% GENERATED by Thesis/tools/make_tables.py -- do not edit by hand.
% Every figure below is read from a CSV written by the pipelines.
\begin{table}[htbp]
\centering
\small
\caption[What each side of the baseline comparison was given]{What each side of the baseline comparison was given. The variable set is identical within every domain, so a single column carries it; the difference, where there is one, is the subsampling rate and therefore the number of fault-free samples. Robotics and industrial are matched exactly. Healthcare is not, and the direction runs against the thesis: the autoencoders are fitted on 7{,}143 fault-free samples where the best causal configuration sees 5{,}000, so the causal margin reported for that domain is measured against a better-resourced opponent instead of a handicapped one.}
\label{tab:baseline_parity}
\begin{tabularx}{\textwidth}{@{}l X r r r r r c@{}}
\toprule
\multirow{2}{*}{Domain} & \multirow{2}{*}{Best causal} & \multirow{2}{*}{Vars} & \multicolumn{2}{c}{Subsample} & \multicolumn{2}{c}{Fault-free samples} & \multirow{2}{*}{Matched} \\
\cmidrule(lr){4-5}\cmidrule(lr){6-7}
 & & & Neural & Causal & Neural & Causal &  \\
\midrule
Robotics & PCMCI, RBF & 244 & 74 & 74 & 274 & 274 & yes \\
Healthcare & GES, Linear & 12 & 7 & 10 & 7143 & 5000 & no \\
Industrial & PCMCI, Poly-3 & 52 & 1 & 1 & 500 & 500 & yes \\
\bottomrule
\end{tabularx}
\end{table}
